\documentclass[11pt,a4paper]{article}
\usepackage[utf8]{inputenc}
\usepackage[T1]{fontenc}
\usepackage[brazil]{babel}
\usepackage{lmodern}
\usepackage[margin=2cm]{geometry}
\usepackage{amsmath,amssymb}
\usepackage{booktabs}
\usepackage{graphicx}
\usepackage{hyperref}
\usepackage{xcolor}
\usepackage{titlesec}
\usepackage{tabularx}
\usepackage{enumitem}
\usepackage{float}
\usepackage{array}
\renewcommand{\tabularxcolumn}[1]{>{\raggedright\arraybackslash}p{#1}}
\usepackage{newunicodechar}
\usepackage[most]{tcolorbox}
\usepackage{etoolbox}

\hypersetup{colorlinks=true, linkcolor=blue!70!black, citecolor=blue!70!black, urlcolor=blue!70!black}

\titleformat{\section}{\large\bfseries}{\thesection.}{0.5em}{}
\titleformat{\subsection}{\normalsize\bfseries}{\thesubsection.}{0.5em}{}
\titlespacing*{\section}{0pt}{1.2ex plus 0.5ex minus .2ex}{0.6ex}
\titlespacing*{\subsection}{0pt}{1.0ex plus 0.4ex minus .2ex}{0.4ex}
\setlength{\parskip}{0.45em}
\setlength{\parindent}{0pt}
\setlist{itemsep=0.15em,topsep=0.3em}

% --- símbolos ---
\newunicodechar{π}{\ensuremath{\pi}}
\newunicodechar{τ}{\ensuremath{\tau}}
\newunicodechar{φ}{\ensuremath{\varphi}}
\newunicodechar{ψ}{\ensuremath{\psi}}
\newunicodechar{Σ}{\ensuremath{\Sigma}}
\newunicodechar{→}{\ensuremath{\rightarrow}}
\newunicodechar{⇒}{\ensuremath{\Rightarrow}}
\newunicodechar{∈}{\ensuremath{\in}}
\newunicodechar{−}{\ensuremath{-}}
\newunicodechar{≈}{\ensuremath{\approx}}
\newunicodechar{≠}{\ensuremath{\neq}}
\newunicodechar{≡}{\ensuremath{\equiv}}
\newunicodechar{≤}{\ensuremath{\leq}}
\newunicodechar{≥}{\ensuremath{\geq}}
\newunicodechar{′}{\ensuremath{'}}
\newunicodechar{‑}{-}
\newunicodechar{✔}{\ensuremath{\checkmark}}
\newunicodechar{✘}{\ensuremath{\times}}

% --- cores e caixas ---
\definecolor{cC}{HTML}{7A4FA3}
\definecolor{cS}{HTML}{2F7D74}
\definecolor{cAcc}{HTML}{2F5D8A}
\newcommand{\metrictag}[2]{{\setlength{\fboxsep}{1.5pt}%
  \ifstrequal{#1}{c}{\colorbox{cC}{\textcolor{white}{\scriptsize\textbf{#2}}}}{%
  \ifstrequal{#1}{s}{\colorbox{cS}{\textcolor{white}{\scriptsize\textbf{#2}}}}{%
  \colorbox{black!60}{\textcolor{white}{\scriptsize\textbf{#2}}}}}}}
\newcommand{\pagechip}[1]{{\setlength{\fboxsep}{1.5pt}\colorbox{cAcc}{\textcolor{white}{\scriptsize\texttt{#1}}}}}
\newcommand{\readtag}[1]{\texttt{\small #1}}
\newcommand{\partbar}[1]{\par\vspace{1.2em}\noindent\colorbox{black!80}{\parbox{\dimexpr\linewidth-2\fboxsep}{\color{white}\bfseries\sffamily\MakeUppercase{#1}}}\par\vspace{0.4em}}
\newtcolorbox{citacard}[1]{breakable, enhanced, colback=white, colframe=black!25, boxrule=0.5pt, arc=2pt,
  left=6pt, right=6pt, top=4pt, bottom=4pt, title={\small #1}, coltitle=black, colbacktitle=black!6, fonttitle=\normalfont}
\newenvironment{origquote}{\par\itshape\small}{\par}
\newenvironment{tradquote}{\par\small\textsc{Tradução:}\ }{\par}
\newcommand{\explab}[1]{\par\smallskip{\small\color{cAcc}\textbf{\MakeUppercase{#1}}}\par}
\colorlet{notebg}{blue!4}\colorlet{notefg}{cAcc}
\colorlet{okbg}{green!5}\colorlet{okfg}{green!45!black}
\colorlet{warnbg}{orange!8}\colorlet{warnfg}{orange!80!black}
\newtcolorbox{infobox}[2][note]{breakable, enhanced, boxrule=0pt, arc=2pt, left=6pt, right=6pt, top=4pt, bottom=4pt,
  colback=#1bg, borderline west={3pt}{0pt}{#1fg}, frame hidden,
  before upper={\ifstrempty{#2}{}{\textbf{#2}\par}}}
\colorlet{pairc}{cC}\colorlet{pairs}{cS}
\newenvironment{paircol}[2]{\begin{tcolorbox}[width=0.485\linewidth, nobeforeafter, box align=top, boxrule=0.5pt,
  colframe=black!25, colback=white, borderline north={3pt}{0pt}{pair#1}, left=5pt, right=5pt,
  before upper={{\color{pair#1}\small\textbf{\MakeUppercase{#2}}}\par}]}{\end{tcolorbox}}

\newunicodechar{ℓ}{\ensuremath{\ell}}
\newunicodechar{Φ}{\ensuremath{\Phi}}
\newunicodechar{‖}{\ensuremath{\|}}
\newunicodechar{ρ}{\ensuremath{\rho}}
\newunicodechar{μ}{\ensuremath{\mu}}
\newunicodechar{∫}{\ensuremath{\int}}
\newunicodechar{∞}{\ensuremath{\infty}}

\begin{document}

\begin{center}
    {\LARGE \textbf{Fichamento: Max-Sensitivity e Avg-Sensitivity}} \\[0.35em]
    {\large YEH, C.-K.; HSIEH, C.-Y.; SUGGALA, A. S.; INOUYE, D. I.; RAVIKUMAR, P. \textit{On the (In)fidelity and Sensitivity of Explanations}. NeurIPS, 2019 (arXiv:1901.09392v4)} \\[0.3em]
    {\normalsize Complementar: BHATT, U.; WELLER, A.; MOURA, J. M. F. \textit{Evaluating and Aggregating Feature-based Model Explanations}. IJCAI, 2020} \\[0.3em]
    \textbf{Métricas no Quantus:} \texttt{quantus.MaxSensitivity} e \texttt{quantus.AvgSensitivity} \\[0.2em]
    \textbf{Aluno:} Arthur Santos de Oliveira (RA: 156297) \\[0.2em]
    \textbf{Atualizado em:} 20/10/2026 \\[0.2em]
    \textbf{Repositório:} \url{https://github.com/Arthur-so/shap-ecg-arrhythmia}
\end{center}
\vspace{-0.5em}
\hrule
\vspace{0.8em}

\textbf{Como ler.} \texttt{[Yeh p. N]} indica a página do PDF de Yeh et al. (25 páginas, incluindo o apêndice) e \texttt{[Bhatt p. N]}, a página do PDF de Bhatt et al. (1 a 7). As citações estão no original em inglês, cada uma seguida de uma tradução livre. A \textbf{Parte I} explica as duas métricas sem pressupor a leitura dos artigos; a \textbf{Parte II} é o fichamento; a \textbf{Parte III} descreve a versão implementada no TCC, sua aderência aos artigos, o diagnóstico e os resultados; a versão antiga, do Quantus, aparece resumida no final.

\begin{infobox}[note]{Quem define o quê}
\begin{itemize}
\item \textbf{Max-Sensitivity:} proposta por \textbf{Yeh et al. (2019)}, Definição 3.1. É a referência principal do Quantus para as duas métricas.
\item \textbf{Avg-Sensitivity:} Yeh et al. definem só o máximo. A versão com a \textbf{média} aparece como métrica em \textbf{Bhatt et al. (2020)}, Definição 2, que também retoma a Max (Definição 1). O Quantus cita os dois artigos.
\item \textbf{Correção:} o fichamento da Faithfulness Correlation diz que o critério de baixa sensibilidade de Bhatt et al. “dá origem” às duas métricas. A formulação mais precisa é: a Max foi proposta por Yeh et al. (2019); Bhatt et al. (2020) a retomam e acrescentam a Avg.
\end{itemize}
\end{infobox}

\partbar{Parte I · Entendendo as métricas}

\section{O problema que as métricas resolvem}
As métricas de fidelidade (Faithfulness Correlation e Estimate) perguntam se a explicação descreve o que o modelo faz. A sensibilidade faz outra pergunta: \textbf{se o ECG mudar um pouquinho, a explicação muda muito?} Se duas entradas quase idênticas, com a mesma predição, recebem explicações muito diferentes, é difícil confiar em qualquer uma delas. Uma explicação \textbf{estável} muda pouco quando a entrada muda pouco. No Quantus, essas métricas ficam na categoria \textbf{robustez}.

\section{A ideia em uma frase e um exemplo}
\begin{infobox}[note]{Definição em uma frase}
Geram-se várias versões levemente perturbadas da entrada; para cada uma, calcula-se a explicação de novo e mede-se a \textbf{distância} entre a explicação nova e a original. A \textbf{Max-Sensitivity} é a maior dessas distâncias; a \textbf{Avg-Sensitivity}, a média. Quanto \textbf{menor}, mais estável a explicação.
\end{infobox}

\textbf{Exemplo com 4 atributos.} Um modelo recebe uma entrada com 4 atributos e, para ela, a explicação original dá as importâncias (0,5; 0,3; 0,1; 0,1). A entrada é então levemente alterada de três formas diferentes (as \textbf{perturbações} 1, 2 e 3, detalhadas na seção 3), e cada versão alterada é explicada de novo. A tabela mostra só as \textbf{explicações}, isto é, as importâncias de cada atributo:
\begin{center}\small
\begin{tabular}{lccccc}
\toprule
& \textbf{At. 1} & \textbf{At. 2} & \textbf{At. 3} & \textbf{At. 4} & \textbf{Distância à original} \\
\midrule
Explicação original & 0,50 & 0,30 & 0,10 & 0,10 & --- \\
Após a perturbação 1 & 0,45 & 0,35 & 0,10 & 0,10 & 0,07 \\
Após a perturbação 2 & 0,50 & 0,30 & 0,15 & 0,05 & 0,07 \\
Após a perturbação 3 & 0,10 & 0,30 & 0,10 & 0,50 & 0,57 \\
\bottomrule
\end{tabular}
\end{center}
A distância é a euclidiana: a raiz da soma dos quadrados das diferenças, atributo a atributo. Nas perturbações 1 e 2, a explicação quase não muda; na 3, ela troca o atributo mais importante. A \textbf{Max-Sensitivity} é 0,57 (o pior caso) e a \textbf{Avg-Sensitivity} é $(0{,}07 + 0{,}07 + 0{,}57)/3 = 0{,}24$. A Max é dominada por um único caso ruim; a Avg diz quanto a explicação varia “tipicamente”.

\section{Passo a passo: o que é perturbar a entrada}
A tabela da seção 2 mostra só as explicações. Falta dizer o que é “alterar levemente a entrada” e quão pequena essa alteração precisa ser. Para isso, é preciso separar duas coisas:
\begin{itemize}
\item a \textbf{entrada}: os valores dos atributos que o modelo recebe (no ECG, o valor do sinal em cada derivação e instante);
\item a \textbf{explicação}: a importância que o método dá a cada atributo (a tabela da seção 2).
\end{itemize}
A perturbação mexe na \textbf{entrada}; a sensibilidade observa o que acontece com a \textbf{explicação}.

\subsection{Perturbação e raio}
No exemplo, a entrada original tem os valores (0,80; $-$0,20; 1,50; 0,30). \textbf{Perturbar} é somar a cada atributo um valor pequeno, para mais ou para menos. Para dizer o que é “pequeno”, fixa-se um limite, o \textbf{raio} $r$: nenhum atributo pode mudar mais que $r$. Com $r = 0{,}1$:
\begin{center}\small
\begin{tabular}{lccccc}
\toprule
\textbf{Entrada} & \textbf{At. 1} & \textbf{At. 2} & \textbf{At. 3} & \textbf{At. 4} & \textbf{Maior mudança} \\
\midrule
Original & 0,80 & $-$0,20 & 1,50 & 0,30 & --- \\
Perturbação 1 & 0,86 & $-$0,27 & 1,46 & 0,39 & 0,09 \\
Perturbação 2 & 0,75 & $-$0,11 & 1,53 & 0,28 & 0,09 \\
Perturbação 3 & 0,77 & $-$0,15 & 1,58 & 0,24 & 0,08 \\
\midrule
Não permitida & 0,81 & $-$0,22 & 1,75 & 0,30 & 0,25 \\
\bottomrule
\end{tabular}
\end{center}
Nas perturbações 1 a 3, cada atributo mudou no máximo 0,09, dentro do raio. A última linha não é permitida: o atributo 3 mudou 0,25, mais que $r = 0{,}1$. Repare que a perturbação 3 é tão pequena quanto as outras na entrada, mas, na seção 2, foi ela que fez a explicação trocar o atributo mais importante. É exatamente esse tipo de caso que a métrica procura.

\subsection{Vizinhança}
A \textbf{vizinhança de raio $r$} é o conjunto de \textbf{todas} as entradas permitidas: todas aquelas em que nenhum atributo difere do original mais que $r$. As perturbações 1 a 3 estão na vizinhança; a última linha está fora. Quanto maior o raio, maior a vizinhança e maiores as alterações permitidas. O nome técnico dessa forma de medir a alteração, pela maior mudança entre os atributos (a última coluna da tabela), é \textbf{norma $\ell_\infty$}; por isso os artigos falam em “vizinhança $\ell_\infty$ de raio $r$”.

\subsection{Os passos}
\begin{enumerate}
\item \textbf{Escolher o raio $r$}: o quanto cada atributo da entrada pode mudar. Isso define a vizinhança.
\item \textbf{Gerar entradas dentro da vizinhança.} Em Yeh et al., para cada atributo sorteia-se um valor entre $-r$ e $+r$ e soma-se ao valor original, como nas perturbações 1 a 3. Em Bhatt et al., em vez de sortear, usam-se \textbf{exemplos reais do treino} que estejam dentro da vizinhança e tenham a \textbf{mesma predição} da entrada original.
\item \textbf{Explicar cada entrada gerada} com o mesmo método usado na original.
\item \textbf{Medir a distância} entre cada explicação nova e a original (a última coluna da tabela da seção 2).
\item \textbf{Resumir:} a maior distância é a Max-Sensitivity; a média, a Avg-Sensitivity.
\end{enumerate}

\section{Dois cuidados de leitura}
\begin{itemize}
\item \textbf{Estabilidade não é qualidade por si só.} Uma explicação que devolve sempre o mesmo mapa, qualquer que seja o ECG, tem sensibilidade zero e não explica nada. Os próprios autores dizem isso (Parte II). A sensibilidade só faz sentido lida ao lado da fidelidade.
\item \textbf{Parte da variação pode vir do modelo.} Se o próprio modelo muda muito com a perturbação, é esperado que a explicação também mude. A sensibilidade alta, nesse caso, reflete o modelo, e não um defeito da explicação.
\end{itemize}

\section{Traduzindo para o ECG}
No TCC, $x$ é um ECG de 12 derivações × 15\,000 instantes (escala z-score por derivação), e a explicação é o mapa GradientSHAP com 180\,000 valores. Cada ponto do sinal (uma derivação em um instante) é um atributo da entrada. Perturbar o ECG é somar a cada ponto um valor sorteado entre $-r$ e $+r$, como na seção 3, o que gera um ECG dentro da vizinhança de raio $r$; “explicar de novo” é rodar o GradientSHAP outra vez nesse ECG perturbado. A distância é calculada entre o mapa novo e o original.

\partbar{Parte II · Fichamento dos artigos}

\section{Mapa de leitura}
\begin{center}\small
\begin{tabularx}{\linewidth}{llll}
\toprule
\textbf{Páginas} & \textbf{Seção} & \textbf{Prioridade} & \textbf{Por quê} \\
\midrule
Yeh p. 1--2 & Resumo, §1 & \readtag{ler} & Motivação da sensibilidade e da explicação constante \\
Yeh p. 2--5 & §2 & \readtag{consulta} & Infidelity (métrica a revisar depois) \\
Yeh p. 5 & §3 & \readtag{ler} & Definição 3.1 (Max-Sensitivity) e objeções \\
Yeh p. 6 & §4 & \readtag{consulta} & Suavização reduz a sensibilidade \\
Yeh p. 6--7 & §5, Tabela 1 & \readtag{ler} & Experimentos e valores de referência \\
Yeh p. 16 & Apêndice B.1 & \readtag{ler} & Raio, norma, nº de amostras, normalização \\
Bhatt p. 2 & §3.1 & \readtag{ler} & Definições 1 e 2 (Max e Average Sensitivity) \\
Bhatt p. 5--6 & §7, §7.2 & \readtag{ler} & Distâncias, vizinhança com exemplos do treino \\
\bottomrule
\end{tabularx}
\end{center}

\section{Citações comentadas: Yeh et al. (2019)}

\begin{citacard}{\pagechip{Yeh p. 2} §1 Introdução}
\begin{origquote}“[...] the sensitivity of an explanation, which measures the degree to which the explanation is affected by insignificant perturbations from the test point. It is natural to wish for our explanation to have low sensitivity, since that would entail differing explanations with minor variations in the input (or prediction values), which might lead us to distrust the explanations.”\end{origquote}
\begin{tradquote}“[...] a sensibilidade de uma explicação, que mede o quanto a explicação é afetada por perturbações insignificantes do ponto de teste. É natural desejar que a explicação tenha baixa sensibilidade, pois o contrário implicaria explicações diferentes para variações mínimas na entrada (ou nos valores da predição), o que poderia nos levar a desconfiar das explicações.”\end{tradquote}
\explab{Explicação}
A palavra-chave é \textbf{insignificantes}: a perturbação deve ser pequena o bastante para que a entrada continue, na prática, a mesma. No ECG, isso exige escolher um ruído que não altere o conteúdo clínico do sinal.
\end{citacard}

\begin{citacard}{\pagechip{Yeh p. 2} §1 Introdução}
\begin{origquote}“[...] the optimal explanation that minimizes sensitivity (for all its related variants) is simply a trivial constant explanation, which is naturally not a desired explanation. So a more interesting question would be: how do we modify a given explanation so that it has lower sensitivity, but not too much. To quantify the latter, we could in turn use fidelity.”\end{origquote}
\begin{tradquote}“[...] a explicação ótima que minimiza a sensibilidade (em todas as suas variantes) é simplesmente uma explicação constante trivial, que naturalmente não é desejada. Assim, uma pergunta mais interessante seria: como modificar uma explicação para que tenha menor sensibilidade, mas não demais. Para quantificar esse ‘não demais’, podemos usar a fidelidade.”\end{tradquote}
\explab{Explicação}
Os autores dizem com todas as letras que a sensibilidade sozinha não basta: ela deve ser lida junto com a fidelidade. No TCC, isso significa interpretar a sensibilidade de cada classe ao lado da Faithfulness Correlation e da Faithfulness Estimate.
\end{citacard}

\begin{citacard}{\pagechip{Yeh p. 5} §3 Definição 3.1}
\begin{origquote}“[...] we consider a closely related measure, we term max-sensitivity, that measures the maximum change in the explanation with a small perturbation of the input x. Definition 3.1. Given a black-box function f, explanation functional Φ, and a given input neighborhood radius r, we define the max-sensitivity for explanation as: SENS\textsubscript{MAX}(Φ, f, x, r) = max\textsubscript{‖y−x‖≤r} ‖Φ(f, y) − Φ(f, x)‖.”\end{origquote}
\begin{tradquote}“[...] consideramos uma medida próxima, que chamamos de max-sensitivity, que mede a maior mudança na explicação diante de uma pequena perturbação da entrada x. Definição 3.1. Dados uma função caixa-preta f, um funcional de explicação Φ e um raio de vizinhança r, definimos a max-sensitivity da explicação como: [fórmula acima].”\end{tradquote}
\explab{Explicação}
$\Phi(f, y)$ é a explicação da entrada perturbada $y$, e $\Phi(f, x)$ a da original. A métrica é a maior distância entre as duas, entre todas as entradas $y$ a no máximo $r$ de $x$. A definição não divide pela norma da explicação original: a métrica é uma diferença \textbf{absoluta}.
\end{citacard}

\begin{citacard}{\pagechip{Yeh p. 5} §3 Estimativa por Monte Carlo}
\begin{origquote}“The main attraction of the max-sensitivity measure is that it can be robustly estimated via Monte-Carlo sampling, as in our experiments.”\end{origquote}
\begin{tradquote}“A principal vantagem da max-sensitivity é que ela pode ser estimada de forma robusta por amostragem de Monte Carlo, como nos nossos experimentos.”\end{tradquote}
\explab{Explicação}
O máximo verdadeiro exigiria testar todas as entradas da vizinhança. Na prática, sorteiam-se algumas perturbações e toma-se o maior valor observado. Com poucas amostras, o valor estimado tende a ficar abaixo do máximo verdadeiro.
\end{citacard}

\begin{citacard}{\pagechip{Yeh p. 5} §3 Duas objeções}
\begin{origquote}“[...] sensitivity provides only a partial measure of what is desired from an explanation. [...] The second objection is that natural explanations might have a certain amount of sensitivity by their very nature, either because the model is sensitive, or because the explanations themselves are constructed by measuring the sensitivities of the predictor function, so that their sensitivities in turn is likely to be more than that of the function.”\end{origquote}
\begin{tradquote}“[...] a sensibilidade é apenas uma medida parcial do que se deseja de uma explicação. [...] A segunda objeção é que explicações naturais podem ter certa sensibilidade por sua própria natureza, seja porque o modelo é sensível, seja porque as próprias explicações são construídas medindo a sensibilidade da função preditora, de modo que a sensibilidade delas tende a ser maior que a da função.”\end{tradquote}
\explab{Explicação}
O GradientSHAP é construído a partir de gradientes, ou seja, da sensibilidade do modelo. Por isso, é esperado que ele tenha alguma sensibilidade; o valor isolado não diz se ela é “alta demais”. Comparar classes, ou comparar com um valor de referência, é mais informativo que o número absoluto.
\end{citacard}

\begin{citacard}{\pagechip{Yeh p. 6} §4 Suavização}
\begin{origquote}“Given any kernel k(x, z) [...] we can define a smoothed explanation as Φ\textsubscript{k}(f, x) := ∫ Φ(f, z) k(x, z) dz. When k(x, z) is set to the Gaussian kernel, Φ\textsubscript{k}(f, x) becomes Smooth-Grad. We now show that the smoothed explanation is less sensitive than the original sensitivity averaged around x.”\end{origquote}
\begin{tradquote}“Dado qualquer núcleo k(x, z) [...] podemos definir uma explicação suavizada como [fórmula]. Quando k(x, z) é o núcleo gaussiano, a explicação suavizada é o Smooth-Grad. Mostramos agora que a explicação suavizada é menos sensível que a sensibilidade original média ao redor de x.”\end{tradquote}
\explab{Explicação}
Fazer a média de explicações calculadas em vários pontos vizinhos reduz a sensibilidade. O GradientSHAP faz exatamente isso: soma um ruído gaussiano à entrada e faz a média sobre várias amostras (\texttt{n\_samples}) e baselines. Logo, o número de amostras usado no GradientSHAP afeta diretamente a sensibilidade medida (Parte III).
\end{citacard}

\begin{citacard}{\pagechip{Yeh p. 7} Tabela 1}
\begin{origquote}SENS\textsubscript{MAX} das explicações globais: MNIST — Grad 0,56; IG 0,47; SHAP 0,35; Square 0,24. Cifar-10 — Grad 1,15; IG 1,08; SHAP 0,93. Imagenet — Grad 1,16; IG 0,93.\end{origquote}
\begin{tradquote}(Valores da tabela; não há texto a traduzir.)\end{tradquote}
\explab{Explicação}
Como as explicações são normalizadas para ter norma 1 antes do cálculo (apêndice), a distância entre duas delas vai de 0 (iguais) a 2 (opostas). Os valores servem como ordem de grandeza: 0,35 para o KernelSHAP no MNIST, e cerca de 1 em imagens coloridas. Note que o SHAP do artigo é o KernelSHAP, e não o GradientSHAP.
\end{citacard}

\begin{citacard}{\pagechip{Yeh p. 16} Apêndice B.1 Configuração}
\begin{origquote}“To compute the sensitivity scores defined in Definitions 3.1, we randomly sample 50 points with Monte-Carlo sampling. Following the adversarial literature, we set the norm for ‖y − x‖ [...] to be L∞ norm and the value of the maximum perturbation r is set to 0.1 for all datasets. To allow fair comparisons among different explanation methods, we normalize the explanation to have unit norm before calculating the sensitivity [...].”\end{origquote}
\begin{tradquote}“Para calcular a sensibilidade da Definição 3.1, sorteamos 50 pontos por Monte Carlo. Seguindo a literatura de ataques adversariais, usamos a norma L∞ para ‖y − x‖ [...] e o valor da perturbação máxima r é 0,1 para todos os conjuntos de dados. Para permitir comparações justas entre métodos de explicação, normalizamos a explicação para ter norma unitária antes de calcular a sensibilidade [...].”\end{tradquote}
\explab{Explicação}
Três escolhas concretas: \textbf{(i)} 50 perturbações por exemplo; \textbf{(ii)} norma L∞ com $r = 0{,}1$, isto é, cada pixel pode mudar no máximo 0,1 (o apêndice não informa a escala das imagens; se for de 0 a 1, como é usual, isso é 10\% da faixa); \textbf{(iii)} as duas explicações são normalizadas para norma 1 antes de medir a distância. A normalização faz a métrica olhar o \textbf{formato} do mapa, e não a sua escala.
\end{citacard}

\section{Citações comentadas: Bhatt et al. (2020)}

\begin{citacard}{\pagechip{Bhatt p. 2} §3.1 Desideratum: Low Sensitivity}
\begin{origquote}“We want to ensure that, if inputs are near each other and their model outputs are similar, then their explanations should be close to each other. [...] Let N\textsubscript{r} = \{z ∈ D\textsubscript{x} | ρ(x, z) ≤ r, f(x) = f(z)\} be a neighborhood of datapoints within a radius r of x.”\end{origquote}
\begin{tradquote}“Queremos garantir que, se as entradas estão próximas e as saídas do modelo são parecidas, então as explicações também estejam próximas. [...] Seja N\textsubscript{r} = [...] uma vizinhança de pontos de dados a um raio r de x.”\end{tradquote}
\explab{Explicação}
Duas diferenças em relação a Yeh et al.: a vizinhança é formada por \textbf{pontos de dados} ($D_x$, exemplos reais), e não por perturbações sorteadas; e exige-se a \textbf{mesma predição}, $f(x) = f(z)$. A segunda condição evita cobrar estabilidade da explicação quando a decisão do próprio modelo mudou.
\end{citacard}

\begin{citacard}{\pagechip{Bhatt p. 2} Definições 1 e 2}
\begin{origquote}“Definition 1 (Max Sensitivity). [...] μ\textsubscript{M}(f, g, r; x) = max\textsubscript{z∈N\textsubscript{r}} D(g(f, x), g(f, z)). Definition 2 (Average Sensitivity). Given [...] a distribution P\textsubscript{x}(·) over the inputs centered at point x, we define the average sensitivity of g at x as: μ\textsubscript{A}(f, g, r; x) = ∫\textsubscript{z∈N\textsubscript{r}} D(g(f, x), g(f, z)) P\textsubscript{x}(z) dz.”\end{origquote}
\begin{tradquote}“Definição 1 (Max Sensitivity). [fórmula]. Definição 2 (Average Sensitivity). Dada [...] uma distribuição P\textsubscript{x}(·) sobre as entradas, centrada em x, definimos a sensibilidade média de g em x como: [fórmula].”\end{tradquote}
\explab{Explicação}
A Average Sensitivity é a média ponderada da distância sobre a vizinhança. Estimada por amostragem, ela é a média das distâncias, e a Max é o máximo das mesmas distâncias. É por isso que o Quantus implementa as duas com o mesmo código, mudando só a última linha.
\end{citacard}

\begin{citacard}{\pagechip{Bhatt p. 5} §7 Configuração}
\begin{origquote}“Before doing aggregation, we unit norm all explanations. [...] We set D = ℓ2 and ρ = ℓ∞.”\end{origquote}
\begin{tradquote}“Antes da agregação, normalizamos todas as explicações para norma 1. [...] Usamos D = ℓ2 e ρ = ℓ∞.”\end{tradquote}
\explab{Explicação}
Mesmas escolhas de Yeh et al.: distância euclidiana ($\ell_2$) entre explicações normalizadas e vizinhança em $\ell_\infty$.
\end{citacard}

\begin{citacard}{\pagechip{Bhatt p. 5--6} §7.2 Max and Avg Sensitivity}
\begin{origquote}“To evaluate the sensitivity criterion, we sample a set of test points from D and an additional larger set of training points. We then find the training points that fall within a radius r neighborhood of each test point and find the distance between each nearby training point explanation and the test point explanation to get a mean and max. [...] Sensitivity is highly dependent on the dimensionality d and on the radius r.”\end{origquote}
\begin{tradquote}“Para avaliar o critério de sensibilidade, sorteamos um conjunto de pontos de teste e um conjunto maior de pontos de treino. Encontramos os pontos de treino que estão a um raio r de cada ponto de teste e calculamos a distância entre a explicação de cada ponto de treino próximo e a do ponto de teste, obtendo a média e o máximo. [...] A sensibilidade depende muito da dimensionalidade d e do raio r.”\end{tradquote}
\explab{Explicação}
Em dados tabulares com poucos atributos, há exemplos reais próximos o bastante. Com 180\,000 valores por ECG, dificilmente existe outro exame do treino a uma distância pequena em norma $\ell_\infty$ (ponto a ponto), então a vizinhança de Bhatt et al. tenderia a ficar vazia. Por isso a versão por perturbações sorteadas, de Yeh et al., é a aplicável ao ECG.
\end{citacard}

\partbar{Parte III · Aplicação ao TCC}

\section{Versão implementada no TCC}
Código: \texttt{src/sensitivity.py} (commit \texttt{ecc19ce}). Para cada um dos 622 mapas do modelo B (ECG × classe explicada, predições corretas do teste):
\begin{enumerate}
\item \textbf{Explicação de referência.} O mapa GradientSHAP do ECG é refeito com a mesma configuração usada para gerar os mapas do TCC (\texttt{src/gshap.py}: referência balanceada de 36 ECGs, 200 sorteios, desvio do ruído 0,09; na versão 1 dos mapas, 32 ECGs sorteados do treino e 50 sorteios) e com uma \textbf{semente aleatória fixa} para aquele mapa.
\item \textbf{Perturbações.} Uma perturbação é uma versão do ECG com um ruído pequeno somado. O ruído é sorteado \textbf{ponto a ponto}: cada ponto de cada derivação do \textbf{trecho real} recebe seu próprio valor, sorteado de forma uniforme entre $-0{,}1$ e $+0{,}1$ (em escala z-score; qualquer valor nesse intervalo tem a mesma chance, então muitos pontos mudam bem menos que 0,1). Não é um deslocamento do sinal inteiro: pontos vizinhos podem subir ou descer. Com quatro pontos, por exemplo:
\begin{center}\small
\begin{tabular}{lrrrr}
\toprule
& \textbf{Ponto 1} & \textbf{Ponto 2} & \textbf{Ponto 3} & \textbf{Ponto 4} \\
\midrule
Original & 0,80 & $-$0,20 & 1,50 & 0,30 \\
\addlinespace
Ruído do sorteio 1 & $+$0,07 & $-$0,03 & $-$0,09 & $+$0,01 \\
Perturbação 1 & 0,87 & $-$0,23 & 1,41 & 0,31 \\
\addlinespace
Ruído do sorteio 2 & $+$0,02 & $+$0,08 & $-$0,04 & $-$0,06 \\
Perturbação 2 & 0,82 & $-$0,12 & 1,46 & 0,24 \\
\bottomrule
\end{tabular}
\end{center}
No ECG real, são 12 derivações com milhares de pontos cada, todos com sorteios próprios. As \textbf{10 perturbações} são 10 sorteios independentes desse ruído: 10 versões levemente diferentes do mesmo ECG, todas a no máximo 0,1 do original em qualquer ponto. Isso é sortear 10 pontos numa vizinhança $\ell_\infty$ de raio $r = 0{,}1$, como em Yeh et al. O preenchimento não é tocado.
\item \textbf{Mesma predição.} Se a perturbação tira a classe explicada da predição (probabilidade abaixo do limiar da classe), ela é descartada, como em Bhatt et al.
\item \textbf{Explicação do ECG perturbado} com a \textbf{mesma semente} da referência: os sorteios internos do GradientSHAP (baselines e ruído) são os mesmos, e a diferença entre os dois mapas reflete só a perturbação da entrada.
\item \textbf{Distância} $\ell_2$ entre as duas explicações normalizadas para norma 1 (Yeh et al., Bhatt et al.). É calculada nas \textbf{36 regiões do batimento} (principal, como nas demais métricas) e ponto a ponto (referência). Vai de 0 (iguais) a 2 (opostas).
\item \textbf{Max-Sensitivity} = maior distância entre as perturbações válidas; \textbf{Avg-Sensitivity} = média. Por classe, reporta-se a média sobre os mapas.
\item \textbf{Piso:} distância entre a referência e outra explicação do mesmo ECG, \textbf{sem perturbação} e com outra semente. Mede quanto o próprio GradientSHAP varia de uma execução para outra.
\end{enumerate}

\section{Aderência aos artigos}
\begin{center}\small
\begin{tabularx}{\linewidth}{lXXl}
\toprule
\textbf{Aspecto} & \textbf{Artigos} & \textbf{Versão do TCC} & \textbf{Aderente?} \\
\midrule
Medida & Maior (Yeh, Def.~3.1) ou média (Bhatt, Def.~2) mudança da explicação na vizinhança & Max e Avg das mesmas perturbações & Sim \\
\addlinespace
Vizinhança & Perturbações sorteadas em $\ell_\infty$ (Yeh, p.~16) & Ruído uniforme por ponto, só no trecho real & Sim \\
\addlinespace
Raio & $r = 0{,}1$ (Yeh, p.~16) & $r = 0{,}1$ em z-score & Sim, adaptado \\
\addlinespace
Mesma predição & Exigida por Bhatt (p.~2); não por Yeh & Perturbações que mudam a predição são descartadas & Sim (Bhatt) \\
\addlinespace
Perturbações por mapa & 50 (Yeh, p.~16) & 10 & \textbf{Parcial} (custo) \\
\addlinespace
Distância & $\ell_2$ entre explicações com norma 1 (Yeh, p.~16; Bhatt, p.~5) & Igual, por regiões e por pontos & Sim, adaptado \\
\addlinespace
Explicação refeita & Mesmo método de explicação & Mesmo método, configuração e semente & Sim (semente: escolha do TCC) \\
\bottomrule
\end{tabularx}
\end{center}

\textbf{Raio.} Yeh et al. usam $r = 0{,}1$ em imagens; no ECG em escala z-score, o mesmo valor é pequeno: cerca de 20\% da amplitude mediana do sinal, e da mesma ordem do deslocamento de linha de base usado no aumento de dados do treino ($\pm$0,1). No diagnóstico, nenhuma perturbação com raio até 0,2 mudou a predição, e o \textit{logit} variou em média 0,13 com $r = 0{,}1$.

\textbf{Semente fixa.} O GradientSHAP é um estimador aleatório: ele sorteia baselines e ruído. Como mostra o diagnóstico a seguir, essa aleatoriedade sozinha muda o mapa muito mais que a perturbação. Usar os mesmos sorteios nos dois lados (uma técnica padrão de redução de variância) isola o efeito da perturbação, que é o que a definição pede. A variação entre sementes não é escondida: ela é reportada como \textbf{piso}.

\textbf{Número de perturbações.} A conta parte da unidade avaliada, o \textbf{mapa}: um par (ECG, classe explicada). São 622 mapas, os mesmos das demais métricas. Para cada mapa, o GradientSHAP é rodado de novo várias vezes:
\begin{itemize}
\item 1 vez no ECG original, com a semente $s$: a \textbf{referência};
\item 1 vez no ECG original, com a semente $s+1$: o \textbf{piso};
\item 1 vez em cada ECG perturbado, com a semente $s$: com 10 perturbações, 10 execuções.
\end{itemize}
São, portanto, $1 + 1 + 10 = 12$ execuções por mapa e $12 \times 622 = 7\,464$ no total, cerca de 7\,500. (Cada execução usa internamente 50 amostras do GradientSHAP; esse 50 é outro parâmetro e não entra na conta. Uma perturbação descartada não chega a ser explicada; na rodada houve só 1.)

Com 50 perturbações por mapa, como em Yeh et al., seriam $1 + 1 + 50 = 52$ execuções por mapa e $52 \times 622 = 32\,344$ no total, cerca de 32\,000: mais de quatro vezes o custo da rodada com 10, que levou 83 min na GPU T4. Por isso foram usadas 10. O preço é que, com menos perturbações, a Max tende a ficar um pouco abaixo do máximo verdadeiro (sorteando menos pontos da vizinhança, é menos provável encontrar o pior caso); a Avg, por ser uma média, é menos afetada.

\section{Diagnóstico}
Com 18 ECGs (2 por classe), na CPU, foram comparados pares de mapas em várias condições:
\begin{center}\small
\begin{tabularx}{\linewidth}{Xrrr}
\toprule
\textbf{Comparação} & \textbf{Regiões} & \textbf{Pontos} & \textbf{Mesmo par top-2} \\
\midrule
Mesmo ECG, outra semente (piso) & 0,37 & 0,89 & 67\% \\
Mapa salvo × refeito (outra semente) & 0,34 & 0,87 & 56\% \\
20 × 50 amostras do GradientSHAP & 0,47 & 1,03 & 50\% \\
\addlinespace
Ruído $r = 0{,}02$, mesma semente & 0,03 & 0,14 & 92\% \\
Ruído $r = 0{,}05$, mesma semente & 0,04 & 0,23 & 92\% \\
Ruído $r = 0{,}1$, mesma semente & \textbf{0,06} & \textbf{0,33} & 97\% \\
Ruído $r = 0{,}2$, mesma semente & 0,10 & 0,49 & 92\% \\
\addlinespace
Configuração antiga do Quantus & 0,43 & 1,02 & 56\% \\
\bottomrule
\end{tabularx}
\end{center}
“Mesmo par top-2” é a fração de casos em que as 2 regiões mais importantes (com sinal) são as mesmas nos dois mapas, que é a explicação usada pela Consistency.

\begin{itemize}
\item \textbf{A explicação é estável diante da perturbação.} Com a mesma semente, um ruído de raio 0,1 muda o mapa por regiões em 0,06 (numa escala até 2) e mantém o par das 2 regiões principais em 97\% dos casos. A sensibilidade cresce de forma suave com o raio.
\item \textbf{A aleatoriedade do GradientSHAP é muito maior que o efeito da perturbação.} Duas execuções no mesmo ECG, sem perturbação nenhuma, diferem em 0,37 por regiões e 0,89 ponto a ponto: cerca de 6 vezes o efeito do ruído de raio 0,1. Em todos os 18 ECGs, o piso foi maior que a maior distância causada pela perturbação, mesmo com $r = 0{,}2$.
\item \textbf{A configuração antiga do Quantus media essa aleatoriedade.} Ela comparava o mapa salvo com um mapa refeito com 20 amostras e outra semente. O valor (0,43 por regiões; 1,28 na distância relativa do Quantus) é praticamente o da linha “20 × 50 amostras”, sem perturbação. Ou seja, a métrica antiga não media a robustez da explicação, e sim o ruído do estimador.
\end{itemize}

\begin{infobox}[warn]{Implicação para as demais métricas}
Os mapas usados em todas as métricas foram gerados com 50 amostras do GradientSHAP, e o diagnóstico mostra que eles têm um ruído de Monte Carlo considerável: duas execuções no mesmo ECG escolhem o mesmo par de regiões principais em só 67\% dos casos. Isso afeta sobretudo a \textbf{Consistency}, que compara exatamente esse par entre ECGs: parte da falta de consistência medida pode vir do ruído do estimador, e não do modelo. As métricas de fidelidade usam os mesmos mapas; o efeito do ruído sobre elas ainda não foi medido. Gerar os mapas com mais amostras reduziria o ruído, ao custo de mais tempo de GPU; é um ponto a avaliar.
\end{infobox}

\section{Como a sensibilidade se relaciona com as outras métricas}
\subsection{Três perguntas diferentes}
As métricas já calculadas no TCC formam duas dimensões: a \textbf{fidelidade} (Faithfulness Correlation e Estimate, que medem a mesma coisa de duas formas e concordam entre si) e a \textbf{consistência} (Consistency). A sensibilidade não é uma terceira medida de nenhuma delas: é uma \textbf{terceira dimensão}, a \textbf{robustez}.
\begin{center}\small
\begin{tabularx}{\linewidth}{lXXl}
\toprule
\textbf{Dimensão} & \textbf{Pergunta} & \textbf{O que se compara} & \textbf{Melhor} \\
\midrule
Fidelidade (FaithCorr, FE) & A explicação descreve o que o modelo faz? & A explicação com a \textbf{saída do modelo}, ao retirar regiões do ECG & Maior \\
\addlinespace
Consistência (Consistency) & ECGs diferentes da mesma classe recebem a mesma explicação? & Explicações de \textbf{ECGs diferentes} & Maior \\
\addlinespace
Robustez (Max/Avg-Sensitivity) & A explicação de um ECG se mantém se o ECG mudar um pouquinho? & Explicações de \textbf{versões quase idênticas do mesmo ECG} & \textbf{Menor} \\
\bottomrule
\end{tabularx}
\end{center}
\begin{itemize}
\item \textbf{Não é a direção da fidelidade.} A fidelidade olha a saída do modelo; a sensibilidade não olha a saída, só compara explicações entre si. Uma explicação pode ser muito estável e nada fiel: o caso extremo é a explicação constante, que dá sempre o mesmo mapa e tem sensibilidade zero (Yeh, p.~2). E pode ser fiel e frágil, se descrever bem o modelo em cada ECG mas mudar muito com um ruído pequeno.
\item \textbf{Não é a direção da Consistency, embora as duas falem de “estabilidade”.} A Consistency compara ECGs diferentes, que podem ter diferenças grandes entre si; a sensibilidade compara um ECG consigo mesmo, com diferenças mínimas. Uma classe pode ter explicações robustas em cada ECG e, ainda assim, pouco consistentes entre ECGs (cada ECG tem a sua razão, estável).
\item \textbf{Direção oposta na leitura.} Nas outras métricas, maior é melhor; aqui, \textbf{menor} é melhor. Ao comparar classes, a melhor posição na sensibilidade é a de menor valor.
\end{itemize}

\subsection{Como ler as três dimensões juntas}
Os próprios autores dizem que a sensibilidade só faz sentido ao lado da fidelidade (Yeh, p.~2 e 5). Combinando as duas:
\begin{center}\small
\begin{tabularx}{\linewidth}{lXX}
\toprule
& \textbf{Pouco sensível (robusta)} & \textbf{Muito sensível (frágil)} \\
\midrule
\textbf{Fiel} & Caso desejado: descreve o modelo e não muda por detalhes irrelevantes. & Descreve o modelo, mas um ruído pequeno muda a explicação; difícil confiar em um mapa isolado. \\
\addlinespace
\textbf{Pouco fiel} & Estável, mas pode ser genérica (próxima da explicação constante). & Pior caso: nem descreve o modelo nem se mantém. \\
\bottomrule
\end{tabularx}
\end{center}
A Consistency acrescenta a terceira camada: se a explicação, além de fiel e robusta, se repete entre ECGs da mesma classe, o modelo usa um critério estável para aquela classe, e o SHAP o mostra.

\subsection{O que a sensibilidade já acrescenta (diagnóstico)}
\begin{enumerate}
\item \textbf{As diferenças entre classes nas outras métricas não vêm de explicações frágeis.} Um ruído pequeno (raio 0,1) muda o mapa por regiões em só 0,06 e preserva o par das 2 regiões principais em 97\% dos casos. As explicações são robustas à entrada; a fidelidade moderada (cerca de 0,6) e as diferenças de Consistency entre classes não se explicam por instabilidade diante de detalhes do sinal.
\item \textbf{A explicação não é “constante”.} Robustez sozinha poderia indicar uma explicação genérica. Como a fidelidade é positiva em todas as classes, esse não é o caso: as explicações são, ao mesmo tempo, estáveis e informativas sobre o modelo.
\item \textbf{O ruído do estimador é uma camada de confiabilidade para todas as métricas.} O piso (0,37 por regiões; par top-2 igual em só 67\% entre execuções) é muito maior que o efeito da perturbação. Isso não é uma propriedade do modelo, e sim do GradientSHAP com 50 amostras, e afeta os mapas usados em todas as métricas. Em especial, ajuda a interpretar a Consistency: o par de regiões principais é estável diante da entrada (97\%), mas não diante do sorteio do estimador (67\%). Parte da falta de consistência entre ECGs pode ser ruído do SHAP, e não do modelo. (Esse diagnóstico usa os mapas da versão 1; com 200 sorteios, o piso cai para 0,18, e a seção 13 mostra que a Consistency quase não muda com isso.)
\end{enumerate}

\subsection{Perguntas que os resultados por classe vão responder}
\begin{itemize}
\item As classes com fidelidade mais baixa (AF, SNR, STD) também têm explicações mais sensíveis? Se sim, a fidelidade baixa vem acompanhada de fragilidade; se não, são explicações estáveis que descrevem pior o modelo.
\item A sensibilidade acompanha o F1? Classes que o modelo classifica melhor têm explicações mais robustas?
\item O piso varia entre classes? Uma classe com piso alto tem a Consistency mais penalizada pelo ruído do estimador, o que muda a leitura da Consistency daquela classe.
\end{itemize}

\section{Resultados}
\def\UrlBreaks{\do\/\do-}\sloppy
\textbf{Duas execuções.} A primeira usou os mapas da versão 1: referência de 32 ECGs sorteados do treino e 50 sorteios (Kaggle, GPU T4, 21/10/2026, 83 minutos; resultados em \texttt{results/modeloB\_sens/}). Desde 26/10/2026 o TCC usa a versão 2 dos mapas, a opção D do documento \url{fichamentos/verificacao-referencia-gradientshap}: referência balanceada, com 4 ECGs de rótulo único de cada uma das 9 classes (36), e 200 sorteios por mapa. A sensibilidade foi refeita com essa configuração (versão 24 do runner, 27/10/2026, 329 minutos, dataset \texttt{arthurso/shap-ecg-modelob-v2}; resultados em \texttt{results/modeloB\_v2/sens/}). Nas duas, foram \textbf{622 mapas} com 10 perturbações de raio 0,1 por mapa, e só 1 das 6\,220 perturbações mudou a predição e foi descartada. As tabelas são da versão 2; a versão 1 aparece para comparação.

\subsection{Por classe}
Distâncias nas 36 regiões (escala de 0 a 2). Classes ordenadas da mais robusta (menor Max-Sensitivity) para a menos robusta. Entre parênteses, a posição da classe em cada medida (na sensibilidade e no piso, 1º = menor valor; nas demais, 1º = maior valor). FE e Consistency também são as da versão 2. A última coluna repete a Max-Sensitivity da versão 1.
\begin{center}\small
\setlength{\tabcolsep}{4pt}
\begin{tabularx}{\linewidth}{Xrrrrrrrr}
\toprule
\textbf{Classe} & \textbf{Mapas} & \textbf{Max} & \textbf{Avg} & \textbf{Piso} & \textbf{FE} & \textbf{Consistency} & \textbf{F1} & \textbf{Max v1} \\
\midrule
BRD  & 172 & 0,029 (1º) & 0,023 & 0,13 (2º) & 0,78 (1º) & 0,72 (1º) & 0,927 (2º) & 0,058 \\
BRE  &  20 & 0,033 (2º) & 0,027 & 0,13 (1º) & 0,70 (2º) & 0,34 (6º) & 0,870 (4º) & 0,047 \\
IAVB &  61 & 0,041 (3º) & 0,032 & 0,18 (5º) & 0,55 (5º) & 0,42 (3º) & 0,884 (3º) & 0,064 \\
STE  &  14 & 0,045 (4º) & 0,036 & 0,18 (4º) & 0,59 (4º) & 0,08 (9º) & 0,596 (9º) & 0,071 \\
CVP  &  59 & 0,047 (5º) & 0,038 & 0,17 (3º) & 0,66 (3º) & 0,38 (5º) & 0,837 (5º) & 0,073 \\
SNR  &  78 & 0,049 (6º) & 0,039 & 0,23 (8º) & 0,49 (8º) & 0,26 (7º) & 0,796 (6º) & 0,077 \\
STD  &  65 & 0,056 (7º) & 0,045 & 0,23 (7º) & 0,49 (9º) & 0,39 (4º) & 0,778 (7º) & 0,085 \\
AF   & 108 & 0,056 (8º) & 0,043 & 0,19 (6º) & 0,52 (7º) & 0,46 (2º) & 0,931 (1º) & 0,081 \\
CAP  &  45 & 0,066 (9º) & 0,054 & 0,25 (9º) & 0,54 (6º) & 0,16 (8º) & 0,750 (8º) & 0,096 \\
\midrule
\textbf{Global} & 622 & \textbf{0,045} & 0,036 & 0,18 & 0,61 & 0,46 & 0,819 & 0,072 \\
\bottomrule
\end{tabularx}
\end{center}
STD e AF empatam na Max-Sensitivity (0,056) e trocam de posição na Avg.

\subsection{O que mudou em relação à versão 1}
\begin{itemize}
\item \textbf{A sensibilidade caiu cerca de 40\%} (Max global de 0,072 para 0,045), em todas as classes. A explicação está no próprio método: os dois mapas comparados (ECG original e perturbado) usam a mesma semente, mas o GradientSHAP avalia o gradiente em pontos diferentes, porque a entrada mudou. Assim, os erros de estimativa dos dois mapas não se cancelam por completo, e uma parte da sensibilidade da versão 1 era esse erro. Com 200 sorteios, o erro diminui e a sensibilidade mede mais a reação da explicação ao ECG.
\item \textbf{O piso caiu pela metade} (0,36 → 0,18), como esperado com quatro vezes mais sorteios. Ele continua maior que a sensibilidade em todos os mapas, agora cerca de 4 vezes (eram 5).
\item \textbf{A ordem das classes quase não mudou:} só BRD e BRE trocam as duas primeiras posições, e AF e STD, a 7ª e a 8ª. Mapa a mapa, a ordem dos valores das duas versões tem correlação de postos de 0,78 (diagnóstico).
\end{itemize}

\subsection{Verificações}
\begin{itemize}
\item \textbf{As explicações são robustas em termos absolutos.} A maior mudança média é 0,045 numa escala até 2, e o desvio dentro de cada classe é pequeno (0,007 a 0,019).
\item \textbf{O ruído do estimador ainda domina em todos os mapas.} Em nenhum dos 622 mapas a Max-Sensitivity passou do piso; em média, o piso é cerca de 4 vezes maior.
\item \textbf{Max e Avg dizem a mesma coisa aqui.} A Max é, em média, 1,26 vez a Avg, e as duas ordenam as classes do mesmo jeito (exceto a troca entre AF e STD). Com perturbações pequenas e uniformes, não há “piores casos” isolados; basta acompanhar uma delas.
\item \textbf{A distância ponto a ponto é dominada por ruído fino do mapa.} Ela é cerca de 5 vezes maior (0,24) e tem um piso de 0,53. Por isso a distância por regiões é a principal, como nas demais métricas.
\end{itemize}

\subsection{Respostas às perguntas da seção 12.4}
\begin{enumerate}
\item \textbf{As classes de fidelidade mais baixa também têm explicações mais sensíveis? Sim.} As 4 classes menos robustas (SNR, STD, AF e CAP, posições 6 a 9) são exatamente as 4 menos fiéis na Faithfulness Estimate (posições 6 a 9), e BRD e BRE são as 2 primeiras nas duas medidas. Mapa a mapa, a relação aparece na mesma direção (correlação de postos de $-$0,52 entre Max-Sensitivity e Faithfulness Estimate, usada só como diagnóstico; $-$0,44 na versão 1). Ou seja, onde o SHAP descreve pior o modelo, ele também é mais frágil.
\item \textbf{A sensibilidade acompanha o F1? Só em parte, como as outras métricas.} BRD (2º no F1, 1º na robustez) e CAP (8º e 9º) coincidem. AF tem o maior F1 e é a 8ª em robustez; STE tem o menor F1 e é a 4ª.
\item \textbf{O piso varia entre classes? Sim, de 0,13 (BRE e BRD) a 0,25 (CAP)}, e isso ajuda a ler a Consistency:
\begin{itemize}
\item \textbf{CAP} tem o maior piso e a 8ª Consistency: parte da baixa consistência de CAP pode vir do ruído do estimador.
\item \textbf{BRE} tem o menor piso e só a 6ª Consistency: a baixa consistência de BRE \textbf{não} é ruído do estimador; as explicações de BRE são estáveis em cada ECG, mas variam entre ECGs (por exemplo, alternando entre o QRS de V1 e o de V2).
\item \textbf{STD} tem piso alto (7º) e mesmo assim Consistency intermediária (4ª): o padrão de STD se repete entre ECGs apesar do ruído.
\end{itemize}
Além disso, cortar o piso pela metade quase não mudou a Consistency: com a referência antiga, passar de 50 para 200 sorteios levou a Consistency global de 0,42 para 0,44 (teste da referência, \texttt{results/modeloB\_basebal}). A falta de consistência entre ECGs vem, portanto, principalmente do modelo, e não do ruído do SHAP, o que corrige a suspeita do item 3 da seção 12.3.
\end{enumerate}

\subsection{O que se conclui}
\begin{enumerate}
\item \textbf{Conceitualmente, a robustez é uma terceira dimensão; nos resultados, ela acompanha a fidelidade.} As classes com explicações mais fiéis são também as mais robustas, nas duas versões dos mapas. Com os dados atuais, as métricas se organizam em dois blocos: \textbf{fidelidade e robustez} de um lado, \textbf{consistência} de outro. A sensibilidade corrobora a fidelidade, em vez de trazer uma ordenação nova.
\item \textbf{Reforça a leitura por tipo de classe.} BRE e BRD, definidas pela forma do QRS, têm explicações fiéis e robustas. AF, SNR e CAP, ligadas ao ritmo ou ao batimento como um todo, e STD, têm explicações menos fiéis e menos robustas.
\item \textbf{Afasta a explicação “frágil” para as diferenças entre classes.} Mesmo nas classes menos robustas, a explicação muda pouco com a perturbação (no máximo 0,07 por regiões).
\item \textbf{O ruído do GradientSHAP continua sendo a principal incerteza de um mapa isolado, mas não decide as conclusões por classe.} Com 200 sorteios, a variação entre execuções caiu pela metade e ainda é cerca de 4 vezes maior que o efeito da perturbação. Mesmo assim, a ordem das classes na sensibilidade, na fidelidade e na Consistency praticamente não mudou entre as versões dos mapas.
\end{enumerate}

\section{Versão antiga (Quantus) e por que foi substituída}
A configuração usada antes (\texttt{quantus.MaxSensitivity} e \texttt{AvgSensitivity}, com 10 amostras) somava ruído de $\pm$0,2 a todos os 180\,000 pontos, inclusive ao preenchimento; refazia a explicação com 20 amostras do GradientSHAP, enquanto os mapas originais tinham 50; não fixava a semente; não descartava perturbações que mudam a predição; e dividia a diferença pelo tamanho do mapa original, em vez de normalizar as duas explicações como no artigo. O diagnóstico acima mostra que o valor resultante era dominado pela aleatoriedade do GradientSHAP. A versão antiga também nunca chegou a produzir resultados no TCC, porque estourava a memória da GPU.

\end{document}
